\subsection{CANONICAL BILINEAR FORM}

We have seen (no. 1, Prop. 3) that the symmetric bilinear form

\[
(x, y) \mapsto \mathrm{B} _ {\mathrm{R}} (x, y) = \sum_ {\alpha \in \mathbb {R}} \langle \alpha^ {\prime}, x \rangle \langle \alpha^ {\prime}, y \rangle
\]

on V is non-degenerate and invariant under A(R). Interchanging the roles of R and \(R^{\prime}\), it follows that the symmetric bilinear form \((x^{*}, y^{*}) \mapsto \mathrm{B}_{\mathbb{R}^{\prime}}(x^{*}, y^{*}) = \sum_{\alpha \in \mathbb{R}^{\prime}} \langle \alpha, x^{*} \rangle \langle \alpha, y^{*} \rangle\) on \(V^{*}\) is non-degenerate and invariant under A(R).

The inverse form of \(B_{R^{\prime}}\) (resp. \(B_{R}\) ) on V (resp. \(V^{*}\) ) will be called the canonical bilinear form on V (resp. \(V^{*}\) ) and denoted by \(\Phi_{R}\) (resp. \(\Phi_{R^{\prime}}\) ). It is non-degenerate and invariant under A(R). Let \(\sigma\) be the isomorphism from V to \(V^{*}\) defined by \(B_{R^{\prime}}\) . Then, for \(x \in V\) and \(y \in V\) :

\[
\Phi_ {R} (x, y) = B _ {R ^ {\prime}} (\sigma (x), \sigma (y)) = \sum_ {\alpha \in R} \langle \alpha , \sigma (x) \rangle \langle \alpha , \sigma (y) \rangle .
\]

But \(\langle \alpha, \sigma(x) \rangle = B_{R^{\prime}}(\sigma(\alpha), \sigma(x)) = \Phi_R(\alpha, x)\) . Hence,

\[
\Phi_ {\mathsf {R}} (x, y) = \sum_ {\alpha \in \mathsf {R}} \Phi_ {\mathsf {R}} (\alpha , x) \Phi (\alpha , y).\tag{16}
\]

In view of Prop. 7 of no. 2, \(\Phi_{R}\) is the only non-zero symmetric bilinear form invariant under W(R) which satisfies the identity (16).

For \(\beta \in \mathbb{R}\) , (16) gives

\[
\varPhi_ {\mathrm{R}} (\beta , \beta) = \sum_ {\alpha \in \mathrm{R}} \varPhi_ {\mathrm{R}} (\alpha , \beta) ^ {2} = \frac {1}{4} \varPhi_ {\mathrm{R}} (\beta , \beta) ^ {2} \sum_ {\alpha \in \mathrm{R}} n (\alpha , \beta) ^ {2},
\]

hence

\[
4 \Phi_ {R} (\beta , \beta) ^ {- 1} = \sum_ {\alpha \in R} n (\alpha , \beta) ^ {2}.\tag{17}
\]

Moreover, by Lemma 2 of no. 1, we have, for \(x, y \in V\) :

\[
\begin{array}{r l} & {\mathrm{B} _ {\mathrm{R}} (x, y) = \sum_ {\alpha \in \mathrm{R}} \varPhi_ {\mathrm{R}} \left(\frac {2 \alpha}{\varPhi_ {\mathrm{R}} (\alpha , \alpha)}, x\right) \left(\frac {2 \alpha}{\varPhi_ {\mathrm{R}} (\alpha , \alpha)}, y\right)} \\ & {\qquad = 4 \sum_ {\alpha \in \mathrm{R}} \varPhi_ {\mathrm{R}} (\alpha , x) \varPhi_ {\mathrm{R}} (\alpha , y) \varPhi_ {\mathrm{R}} (\alpha , \alpha) ^ {- 2}.} \end{array}
\]

It follows that, if \(R\) is irreducible, there exists a constant \(\gamma(R) > 0\) such that

\[
\sum_ {\alpha \in R} \Phi_ {R} (\alpha , x) \Phi_ {R} (\alpha , y) \Phi_ {R} (\alpha , \alpha) ^ {- 2} = \gamma (R) \Phi_ {R} (x, y).\tag{18}
\]

By the definition of \(\gamma (\mathbb{R})\) , we have \(B_{\mathbb{R}}(x,y) = 4\gamma (\mathbb{R})\Phi_{\mathbb{R}}(x,y)\) , so

\[
\Phi_ {\mathrm{R} ^ {\prime}} \left(x ^ {*}, y ^ {*}\right) = (4 \gamma (\mathrm{R})) ^ {- 1} \mathrm{B} _ {\mathrm{R} ^ {\prime}} \left(x ^ {*}, y ^ {*}\right)
\]

for \(x^{*},y^{*}\in V^{*}\) . This proves that \(\gamma (\mathsf{R}) = \gamma (\mathsf{R}^{\prime})\) . On the other hand, for \(\beta \in \mathbb{R}\) ,

\[
\begin{array}{c} \Phi_ {\mathrm{R} ^ {\prime}} (\beta^ {\prime}, \beta^ {\prime}) = (4 \gamma (\mathrm{R})) ^ {- 1} \sum_ {\alpha \in \mathfrak {R}} \langle \beta^ {\prime}, \alpha \rangle^ {2} \\ = \gamma (\mathrm{R}) ^ {- 1} \sum_ {\alpha \in \mathbb {R}} \frac {\Phi_ {\mathrm{R}} (\alpha , \beta) ^ {2}}{\Phi_ {\mathrm{R}} (\beta , \beta) ^ {2}} \end{array}
\]

so, by (16),

\[
\Phi_ {\mathrm{R} ^ {\prime}} (\beta^ {\prime}, \beta^ {\prime}) = \gamma (\mathrm{R}) ^ {- 1} \Phi_ {\mathrm{R}} (\beta , \beta) ^ {- 2} \Phi_ {\mathrm{R}} (\beta , \beta)
\]

or finally

\[
\Phi_ {R} (\beta , \beta) \Phi_ {R ^ {\prime}} (\beta^ {\prime}, \beta^ {\prime}) = \gamma (R) ^ {- 1}.\tag{19}
\]

Further, if all the roots of \(\mathbb{R}\) have the same length \(\lambda\) for \(\varPhi_{\mathbb{R}}\) , (16) and (18) show that

\[
\gamma (R) = \lambda^ {- 4}.\tag{20}
\]

Moreover, if \(h\) is the Coxeter number of \(W\) , the Cor. of Th. 1 of Chap. V, § 6, no. 2 shows that:

\[
h \Phi_ {R} (x, x) = \sum_ {\alpha \in R} \left(x, \frac {\alpha}{\lambda}\right) ^ {2} \text {   for   all   } x \in V.
\]

Comparing with (16), we deduce that

\[
\lambda = h ^ {- 1 / 2} \text { or } \gamma (\mathrm{R}) = h ^ {2}.\tag{21}
\]

Finally, formula (19) shows that the roots of \(R^{\prime}\) have length \(\lambda\) for \(\Phi_{R^{\prime}}\) .

